\begin{proofbox}
  \textbf{\textcolor{CProof}{方法一：对应点法（最直观、最通用）}}
  \textbf{想法：} 不管函数长什么样，只要把原图像上每一个点按要求移动，就能找到新解析式。

  \begin{description}[style=nextline,leftmargin=2.8em,labelsep=0.5em,itemsep=0.4em]
    \item[\textbf{第1步：选一个原来的点。}]
      原图像 $y=f(x)$ 上任取 $P(a,b)$，则 $b=f(a)$。
    \item[\textbf{第2步：看这个点移到哪里。}]
      横坐标变为原来的 $k$ 倍、纵坐标不变后，得到 $P'(ka,b)$。
    \item[\textbf{第3步：用新的横坐标表示原来的。}]
      设新坐标为 $(x,y)$，则 $x=ka$，所以 $a=x/k$；同时 $y=b=f(a)$。
    \item[\textbf{第4步：写出新解析式。}]
      \[
        \boxed{\ y=f\!\left(\frac{x}{k}\right)\ }.
      \]
      由于原来的点可以任意选，所以每一个点都按这个规则移动；反过来，这个式子上的每个点也都能由原图像的点移来。
  \end{description}
  \begin{center}
  \begin{tikzpicture}[x=1.07cm,y=0.85cm,>=stealth,every node/.style={font=\small}]
    \draw[->,gray!75] (-0.2,0)--(4.6,0) node[right] {$x$};
    \draw[->,gray!75] (0,-0.2)--(0,2.8) node[left] {$y$};
    \draw[densely dashed,gray!75] (0,1.6)--(4.1,1.6);
    \fill[blue!75!black] (1.5,1.6) circle (2.3pt);
    \fill[orange!85!black] (3,1.6) circle (2.3pt);
    \draw[densely dashed,blue!75!black] (1.5,0)--(1.5,1.6);
    \draw[densely dashed,orange!85!black] (3,0)--(3,1.6);
    \draw[->,very thick,red!65!black] (1.65,1.88)--(2.85,1.88);
    \node[above left,blue!75!black] at (1.5,1.6) {$P(a,b)$};
    \node[above right,orange!85!black] at (3,1.6) {$P'(ka,b)$};
    \node[below] at (1.5,0) {$a$};
    \node[below] at (3,0) {$ka$};
    \node[left] at (0,1.6) {$b$};
    \node[above] at (2.25,2.12) {以 $k=2$ 为例};
  \end{tikzpicture}

  \end{center}
  \noindent\textbf{图解：} 两个点在同一条水平线上，表示\textbf{函数值没有变}；变化的只是横坐标。

  \medskip
  \textbf{\textcolor{CProof}{方法二：从新坐标倒推（写题最快）}}
  假设横坐标要\textbf{压成原来的一半}。新图像上某个点的横坐标是 $x$，它在原图像上的横坐标必定是 $2x$；由于高度不变，新点的高度就是原函数在 $2x$ 处的高度：
  \[
     y=f(2x).
  \]
  同理，若图像\textbf{拉宽 $3$ 倍}，新横坐标为 $x$ 的点来自原来的 $x/3$，所以
  \[
     y=f\!\left(\frac{x}{3}\right).
  \]
  \begin{center}
  \begin{tikzpicture}[x=0.85cm,y=0.78cm,>=stealth,every node/.style={font=\small}]
    \draw[rounded corners=2pt,fill=blue!5,draw=blue!40] (0,0) rectangle (4.2,1.5);
    \draw[rounded corners=2pt,fill=orange!6,draw=orange!50] (5.7,0) rectangle (9.9,1.5);
    \node[align=center] at (2.1,.75) {原图像取输入 $2x$\\高度 $f(2x)$};
    \node[align=center] at (7.8,.75) {新图像取输入 $x$\\高度仍为 $f(2x)$};
    \draw[->,very thick,red!65!black] (4.3,.75)--(5.6,.75);
    \node[above,align=center] at (4.95,2) {横坐标减半};
  \end{tikzpicture}

  \end{center}
  \noindent\textbf{为什么对？} 原来输入 $2x$ 和新的输入 $x$ 指向的是\textbf{同一个移动前后的点}，所以它们的函数值相等。

  \medskip
  \textbf{\textcolor{CProof}{结论回扣}}
  \textbf{横坐标乘 $k$，括号里的 $x$ 除以 $k$。} 其中 $k>0$；若 $k=1$，所有点都在原处。该规律不要求函数关于 $y$ 轴对称，也不要求函数只增不减。
\end{proofbox}
